\ParallelTitle{放慢脚步，重看熟悉的街道}{いつもの通りを、ゆっくり歩く}

\ParallelSection{notice}{留一点时间观察}{気づくための時間をつくる}

\ParallelText{opening}{一条熟悉的街道，有时只是两项日程之间的一段路。不妨试着走一次，不急着完成任何差事。目的不是寻找引人注目的地标，也不是拍出完美的照片，而是留意那些在匆忙赶路时被忽略的细节。}{いつもの通りが、予定と予定をつなぐだけの道になっていることがあります。一度、用事を済ませるためではなく歩いてみましょう。立派な名所を探したり、完璧な写真を撮ったりするのではなく、急いでいると見落とす細部に目を向けるのが目的です。}

\ParallelText{change}{一扇刷过漆的门、一片树荫，或敞开的窗户里传来的碗碟声，都可以成为观察的起点。无须立刻解释每一个细节。有价值的观察可能先是一个问题；散步时可以暂时保留它，等有空再进一步了解。}{塗り直された扉、木陰、開いた窓から聞こえる食器の音。それだけでも観察のきっかけになります。すべてをその場で説明する必要はありません。問いから始まる気づきもあります。調べる時間ができるまで、その問いを残しておいてもよいのです。}

\ParallelText{prompt}{\begin{quote}一种声音、一种颜色、一个问题：给自己的笔记一个简单的提示。\end{quote}}{\begin{quote}音を一つ、色を一つ、問いを一つ。自分のメモのための簡単な手がかりです。\end{quote}}

\ParallelWideFigure{shared-photo}{footpath.png}{两种语言共用一张步道照片。这张照片用于说明散步的环境，并不是建议路线的地图。}{両方の列で共有する小道の写真。散歩の雰囲気を示すためのもので、提案したルートの地図ではありません。}

\ParallelSection{route}{选一条简单的路线}{無理のない道順を選ぶ}

\ParallelText{route-plan}{起点最好方便到达，终点也不应让你必须匆忙折返。比起排得满满的行程，一小段环线往往更合适。出行时仍需查看当地的实际情况；下图是虚构的示意图，并不是某个真实地点的路线指引。}{行きやすい出発点と、帰りを急がずに済む終点を選びます。欲張った行程より、短い周回路のほうが役立つこともあります。歩く場所の状況は実際に確かめてください。下の図は架空の例で、実在する場所への道案内ではありません。}

\ParallelFigure{route-image}{route-zh-Hans.png}{虚构路线上的四个停留点。数字只标识停留位置，不代表距离、时间或无障碍条件。}{四つの立ち止まる場所を示した架空の道順です。数字は場所の目印であり、距離、所要時間、バリアフリー状況を示すものではありません。}[route-ja.png]

\ParallelText{steps.item-1}{\begin{itemize}\item \phantomsection\label{steps}\hypertarget{steps}{}每次停留只选一个细节，不必试图描述周围的一切。\end{itemize}}{\begin{itemize}\item \ifdefstring{\ParallelMode}{right}{\phantomsection\label{steps}\hypertarget{steps}{}}{}立ち止まるたびに一つの細部を選び、周囲のすべてを説明しようとしないこと。\end{itemize}}

\ParallelText{steps.item-2}{\begin{itemize}\item 留一点没有安排的时间。遇到值得再看一眼的事物，可以多停留一会儿。\end{itemize}}{\begin{itemize}\item 予定を入れない時間も残しましょう。もう一度見たいものがあれば、その場にいられます。\end{itemize}}

\ParallelText{steps.item-3}{\begin{itemize}\item 留意他人的隐私。写下一项观察，并不需要给陌生人拍照。\end{itemize}}{\begin{itemize}\item 周囲の人のプライバシーに配慮しましょう。観察を書き留めるために、見知らぬ人を撮影する必要はありません。\end{itemize}}

\ParallelSection{long-observation}{把一次观察写完整}{一つの観察をじっくり追う}

\ParallelProse{long-prose}{设想在两个不同的下午，沿着同一条短路线散步。第一次经过时，一辆送货车停在商店旁，店门不断开合，还没有看见说话的人，交谈声就已经传到了人行道上。第二次经过时，车走了，门也关着。很容易把第一个下午说成“热闹”，把第二个说成“安静”，但这些词已经把许多细小的观察合并成了判断。更有用的笔记，或许可以先写能够核实的事：车停在哪里，门开了多久，以及转过街角后是否还能听见交谈声。这并不要求跟踪陌生人、记录私人谈话，或把一个社区当成供人围观的景观。观察公共环境的同时，仍然可以尊重他人的隐私。走到下一个路口时，找一个不妨碍通行的位置停下，回头看看刚走过的路。从一个方向看很显眼的标牌，换个方向可能被树遮住；原本看起来平坦的人行道，在墙边的参照下也许显出了一点缓坡。就连一张长椅，也能改变这次散步的体验：不停下来时，它只是途中经过的物件；坐下来后，同一个地方却提供了不同的视线高度、不同的景致，也让人有时间听见先前被脚步声盖住的声音。之后写下这段经历时，不必让每个细节都服务于一个整齐的故事。保留散步时原本存在的不确定性。如果店门关着，就写经过时门是关着的，不要据此推断商店已经停业；如果广场看起来空荡荡，也要记得自己只在那里待了几分钟。再去一次的价值，在于增加一份观察，而不是自动证明第一次的印象有误。可以换个时间再来，沿着相反的方向走这条环线，也可以和一个会留意不同事物的人同行。目的不是收集尽可能多的细节，而是认真看清其中的几个，让熟悉的街道变得更容易描述，同时仍然愿意接受下一次散步带来的新发现。 合上笔记本之前，还可以想一想：一个没有到过现场的读者，需要知道什么？“在树旁边”这样的说法，对自己可能很清楚，对别人却未必有用。不妨补上一个不依赖共同记忆也能理解的位置关系：长椅朝向路口，步道绕到树篱后面，或者店门朝着人行道较窄的一侧打开。这些细节的作用是让场景更清楚，而不是把一段虚构示例写得像经过测绘的地图。描述旁边也可以保留一个尚未回答的问题。声音消失了，是因为窗户关上了，因为自己走远了，还是因为另一种声音把它盖住了？只散步一次，也许无法得到可靠的答案。有时，以这个问题结束，会比给出一个自信的解释更准确，也更能引起继续观察的兴趣。下一次来时，可以从这里重新开始，带着好奇心，而不是带着一个预先选好的结论。给笔记标上日期，区分哪些内容是在现场写下的，哪些是后来补充的。记忆可能会把当时彼此独立的细节联系起来。这个小小的区分，有助于后来的读者了解这段记录是如何形成的。}{別々の日の午後に、同じ短い道を散歩する場面を想像してください。最初の訪問では、店の脇に配達のバンが止まり、ドアが何度も開き、話している人の姿が見える前に会話の声が歩道まで届いています。二度目の訪問では、バンは去り、ドアは閉まっています。最初の午後を「にぎやか」、次を「静か」と呼びたくなるかもしれませんが、その言葉はすでに多くの小さな観察を一つの判断にまとめています。より役立つメモは、実際に確かめられることから始められるでしょう。バンがどこに止まっていたか、ドアがどれほどの間開いていたか、角を曲がった後も会話が聞こえたか、といったことです。そのために見知らぬ人を追いかけたり、私的な会話を録音したり、地域を見世物のように扱ったりする必要はありません。人々のプライバシーを尊重しながら、公共の場所を観察できます。次の交差点では、通行を妨げない所で立ち止まり、今歩いてきた道を振り返ってみましょう。一方からは目立っていた標識が、反対からは木に隠れているかもしれません。平らに見えた歩道も、壁と見比べると緩い傾斜に気づくことがあります。ベンチ一つでも散歩の体験は変わります。立ち止まらなければ、ただ通り過ぎる物ですが、腰を下ろせば、同じ場所で目線の高さも眺めも変わり、足音にかき消されていた音に耳を澄ます時間が生まれます。後でその体験を書くときは、すべての細部を一つのきれいな物語に収めようとしないでください。散歩中にあった不確かさを残しましょう。店が閉まっていたなら、廃業したと決めつけず、自分が通ったときにはドアが閉まっていたと書きます。広場に誰もいないように見えても、そこにいたのはほんの数分だったことを忘れないようにしましょう。もう一度訪れることが役立つのは、別の観察が得られるからであり、最初の印象が誤りだったと自動的に証明されるからではありません。違う時間に来たり、周回路を逆向きに歩いたり、違うことに目を留める人と同じ道を歩いたりしてもよいでしょう。目的は、できるだけ多くの細部を集めることではありません。いくつかの細部に十分な注意を向け、なじみのある通りをより描写しやすくすることです。同時に、次の散歩で変わるかもしれないことにも心を開いておきます。ノートをしまう前に、その場にいなかった読者には何が必要か考えてみてください。「木のそば」という表現は、自分には明確でも、他の人には役立たないことがあります。同じ記憶を持たなくても理解できる位置関係を加えましょう。ベンチは交差点に向いていた、小道は生け垣の裏へと曲がっていた、店のドアは歩道の狭い部分に向かって開いた、などです。こうした細部は場面を明確にするためのもので、架空の描写を測量した地図のように見せるためではありません。描写の脇に、答えの出ていない問いを残してもよいでしょう。音が消えたのは、窓が閉まったからなのか、自分が遠ざかったからなのか、それとも別の音にかき消されたからなのか。一度の散歩だけでは、確かな答えは得られないかもしれません。自信に満ちた説明で締めくくるより、その問いで終える方が正確で、先を知りたくなることもあります。次の訪問は、あらかじめ決めた結論ではなく、好奇心を持ってそこから始められます。メモには日付を付け、その場で書いたことと、後から加えたことを区別しましょう。記憶は、そのときには別々だった細部を結びつけることがあります。この小さな区別が、後の読者に、記録がどのように作られたかを伝える助けになります。}

\ParallelSection{notebook}{带一本小笔记本}{小さなノートに書き留める}

\ParallelText{notebook-prose}{笔记本里可以写片段，不必句句都完整漂亮。先记录实际看到的内容，再补充解释。“长椅是空的”是一项观察；“没有人来这个广场”则是范围大得多的判断。把这两类表述分开，之后写出来的文字会更诚实，也更有意思。}{ノートには、整った文章ではなく断片を書いてもかまいません。解釈を加える前に、実際に気づいたことを記録します。「ベンチは空いていた」は観察ですが、「この広場は誰も使わない」はずっと広い主張です。二つを区別すれば、後で書く文章がより誠実で興味深いものになります。}

\begin{ParallelKeep}

\ParallelText{notes-table.row-0}{\phantomsection\label{notes-table}\hypertarget{notes-table}{}\begin{tabularx}{\linewidth}{@{}>{\RaggedRight\arraybackslash}X>{\RaggedRight\arraybackslash}X@{}}\toprule \textbf{观察} & \textbf{问题}\\ \midrule \end{tabularx}}{\ifdefstring{\ParallelMode}{right}{\phantomsection\label{notes-table}\hypertarget{notes-table}{}}{}\begin{tabularx}{\linewidth}{@{}>{\RaggedRight\arraybackslash}X>{\RaggedRight\arraybackslash}X@{}}\toprule \textbf{観察} & \textbf{問い}\\ \midrule \end{tabularx}}

\ParallelText{notes-table.row-1}{\begin{tabularx}{\linewidth}{@{}>{\RaggedRight\arraybackslash}X>{\RaggedRight\arraybackslash}X@{}}一张空长椅 & 换个时间会有人坐吗？\\ \end{tabularx}}{\begin{tabularx}{\linewidth}{@{}>{\RaggedRight\arraybackslash}X>{\RaggedRight\arraybackslash}X@{}}空いているベンチ & 別の時間には使われている？\\ \end{tabularx}}

\ParallelText{notes-table.row-2}{\begin{tabularx}{\linewidth}{@{}>{\RaggedRight\arraybackslash}X>{\RaggedRight\arraybackslash}X@{}}一扇新刷漆的门 & 它以前是什么颜色？\\ \end{tabularx}}{\begin{tabularx}{\linewidth}{@{}>{\RaggedRight\arraybackslash}X>{\RaggedRight\arraybackslash}X@{}}塗り直された扉 & 前は何色だった？\\ \end{tabularx}}

\ParallelText{notes-table.row-3}{\begin{tabularx}{\linewidth}{@{}>{\RaggedRight\arraybackslash}X>{\RaggedRight\arraybackslash}X@{}}窗户里传来音乐 & 走远一点，听起来会不同吗？\\ \bottomrule \end{tabularx}}{\begin{tabularx}{\linewidth}{@{}>{\RaggedRight\arraybackslash}X>{\RaggedRight\arraybackslash}X@{}}窓から聞こえる音楽 & 離れると聞こえ方は変わる？\\ \bottomrule \end{tabularx}}

\ParallelText{notes-table.caption}{这些虚构的笔记展示了观察与问题的区别，只是示例，并非对某个真实街区的调查结果。}{架空のメモで、観察と問いの違いを示しています。実在する地域についての調査結果ではありません。}\end{ParallelKeep}

\ParallelSection{return}{回来以后，再读一遍这段路}{歩いた道を読み返す}

\ParallelText{return-prose}{回来后，可以选一条笔记，展开成一小段文字。保留观察到的细节，说明它为什么对你重要，并标明还有哪些事情尚不清楚。事实性的问题可以稍后核实，但不必把每次散步都变成一个研究项目。有时候，认真留意过，本身就有价值。}{帰ったら、メモを一つ選んで短い段落にしてみます。観察した細部を残し、自分にとって気になった理由を書き、まだ分からないことも明示します。事実に関する問いは後で確かめられますが、散歩を毎回研究課題にする必要はありません。ただ注意を向けたことに価値がある場合もあります。}

\ParallelText{route-reference}{\ParallelReference{route}{回看路线部分}}{\ParallelReference{route}{道順の節に戻る}}
